\section{Evolução Tecnológica na Odontologia}
\label{sec:evolucao}

A trajetória cronológica e instrumental da odontologia digital pode ser intelectualmente compreendida como uma sucessão de ondas tecnológicas superpostas. Cada inflexão histórica não apenas edificou-se sobre as premissas axiomáticas da precedente, como também expandiu assintoticamente o espectro de intervenção e a previsibilidade da práxis clínica \cite{duggal2026exploring}.

\subsection*{Anos 1990: A Primeira Onda Digital e a Desmaterialização dos Registros}

A transição para a década de 1990 assinalou o limiar da informatização nos ambientes clínicos odontológicos. O marco fundamental desta era foi a introdução e posterior comercialização da radiologia digital, instrumentalizada por sensores de estado sólido (dispositivos de carga acoplada - CCD, e semicondutores de óxido metálico complementar - CMOS) e placas de fósforo fotoestimulável (PSP). Esta tecnologia promoveu a obliteração do moroso processo químico de revelação halogenada, logrando uma drástica supressão da dose de radiação ionizante (em até 80\% comparada aos filmes radiográficos tipo E ou F) e viabilizando, primariamente, o armazenamento eletrônico e a transmissibilidade da imagem médica \cite{khoshfekr2025review}. 

De forma simultânea, softwares semânticos de gestão de banco de dados relacionais começaram a suplantar os arquivamentos físicos. A condensação de prontuários, anamneses e registros iconográficos em servidores centralizados instituiu a base de dados longitudinal do paciente.

Embora ainda rudimentar do ponto de vista do processamento algorítmico avançado, esta primeira onda estipulou dois pressupostos inegociáveis: a conversão do sinal físico (fótons/raios-X) em matriz binária (pixels) no momento da aquisição, e a criação de registros perenes passíveis de mineração computacional \cite{mdpi2025convergence}. Sem o alicerce da desmaterialização imagética, as iterativas complexidades subsequentes teriam sido inviabilizadas.

\subsection*{Anos 2000: O Advento CAD/CAM e a Tomografia Volumétrica}

O romper do novo milênio testemunhou o florescimento de duas inovações de caráter marcadamente disruptivo. A primazia pertence à consolidação sistemática dos complexos CAD/CAM (\textit{Computer-Aided Design / Computer-Aided Manufacturing}) em reabilitação oral. Protagonizado pelo sistema Cerec, este fluxo instituiu a viabilidade da fotogrametria intraoral acoplada a um software de desenho assistido e à fresagem subtrativa de monoblocos cerâmicos (feldspáticos ou de dissilicato de lítio) no escopo de uma única intervenção clínica (\textit{chairside}) \cite{alshadidi2024surface}.

A segunda ruptura de paradigma consubstanciou-se na introdução comercial da tomografia computadorizada de feixe cônico (CBCT). Rompendo as amarras da radiografia cefalométrica e panorâmica convencional 2D, a CBCT facultou a reconstrução tridimensional integral do maciço maxilofacial humano, garantindo imagens volumétricas sem magnificação com dosimetria ordens de magnitude inferior à tomografia helicoidal multislice médica \cite{li2025impacts}.

Os dispositivos de escaneamento intraoral (IOS) de primeira e segunda geração, ainda que limitados por deficiências de processamento e necessidade de contrastes em pó, proveram o alicerce empírico de que as nuvens de pontos poderiam substituir a moldagem elastomérica, seja para substratos restauradores, seja para setup ortodôntico \cite{li2024aligner}. \destaque{A precisão passiva submilimétrica conferiu irrefutabilidade acadêmica e aceitação mercadológica à migração do gesso para o voxel e o polígono.}

\subsection*{Anos 2010: Manufatura Aditiva e Orquestração do Planejamento Virtual}

A terceira onda da odontologia digital notabilizou-se pela sinergia profunda entre a aquisição estereoscópica e a manufatura aditiva. O que antes era circunscrito à prototipagem rápida da engenharia aeroespacial e automotiva — a impressão 3D (estereolitografia e processamento digital de luz) — metamorfoseou-se em pedra angular no consultório. O espectro clínico dilatou-se: modelos dentários com resinas fotopolimerizáveis, guias cirúrgicos implantodônticos com inserção e angulação pré-determinadas, splints oclusais, bases para próteses totais e biótipos ósseos para cirurgia ortognática \cite{techsoc2025precision}.

O planejamento virtual reverso, suportado por suítes de software dedicadas, assumiu o protagonismo das reabilitações extensas. A manipulação tridimensional permitiu a simulação cirúrgica profilática, fundindo dados de superfície (STL/PLY) e dados volumétricos (DICOM), alçando a assertividade cirúrgica a patamares de variância menor que 0.2 mm \cite{xing2024clinical}. 

A disciplina ortodôntica, notavelmente, foi reconfigurada pelos alinhadores bolhados em termoplástico, dependentes absolutos da integração digital ponta a ponta: do estagiamento cinemático milimétrico \textit{in silico} à produção seriada automatizada de dispositivos \cite{li2024aligner}.

\subsection*{Anos 2020: A Revolução Cognitiva da Inteligência Artificial e a Pangeia dos Ecossistemas}

A onda hodierna, ainda em pleno curso de expansão, define-se pela amálgama de três vetores cardinais: a inteligência artificial (IA) discriminativa e generativa, a hiperconectividade em nuvem e a imperiosidade dos ecossistemas de arquitetura aberta (\textit{open-source} e \textit{open-architecture}) \cite{infante2024fmi}. Arquiteturas baseadas em redes neurais convolucionais (CNNs) e \textit{transformers} têm superado a acurácia de operadores humanos treinados na detecção incidental de patologias, segmentação multitecidular em CBCT e estadiamento de doença periodontal \cite{elsaid2025digital}.

\destaque{Nesta década, a interoperabilidade sistêmica cessa de ser um luxo técnico para se tornar uma exigência clínica, suplantando os ecossistemas enclausurados (walled gardens) do passado.} Padrões abertos e infraestruturas escaláveis (como o paradigma FHIR e co-simulações FMI/FMU) asseguram que a aquisição, a modelagem física, a simulação biológica e a manufatura computem dados simultaneamente, forjando assim a iminência dos verdadeiros gêmeos digitais aplicados \cite{infante2025distributed}.
